%% suspicious_authors_section.tex
%% ──────────────────────────────────────────────────────────────────
%% New Results subsection: suspicious-author cross-referencing.
%% Include in main.tex via:  %% suspicious_authors_section.tex
%% ──────────────────────────────────────────────────────────────────
%% New Results subsection: suspicious-author cross-referencing.
%% Include in main.tex via:  %% suspicious_authors_section.tex
%% ──────────────────────────────────────────────────────────────────
%% New Results subsection: suspicious-author cross-referencing.
%% Include in main.tex via:  %% suspicious_authors_section.tex
%% ──────────────────────────────────────────────────────────────────
%% New Results subsection: suspicious-author cross-referencing.
%% Include in main.tex via:  \input{suspicious_authors_section}
%% (place after the Mixing Matrix subsection, before \section{Discussion})
%% ──────────────────────────────────────────────────────────────────

\subsection{Cross-referencing Suspicious Authors}\label{sec:suspicious-authors}

To move from aggregate outlier statistics to individual-level findings,
we cross-reference the 189 flagged outliers against real-world author
identities and publication records.
Each ORCID identifier is resolved to a name via the ORCID public
API,\footnote{\url{https://pub.orcid.org}} and publication portfolios
are reconstructed by joining author\,--\,work associations to the
Crossref works table and resolving journal titles through the ISSN
catalogue.

\subsubsection*{Suspiciousness scoring}

We define a composite \emph{suspiciousness score} as the weighted sum
of population-normalised $z$-scores across six behavioural features:
%
\begin{equation}\label{eq:susp}
  S_i = \sum_{k=1}^{6} w_k \cdot z_{ik},
  \qquad
  z_{ik} = \frac{x_{ik} - \bar{x}_k}{\sigma_k},
\end{equation}
%
where $x_{ik}$ is the raw value of feature~$k$ for author~$i$,
$\bar{x}_k$ and $\sigma_k$ are the population mean and standard
deviation, and the weights $w_k$ are chosen to reflect both
Random-Forest feature importance (Figure~\ref{fig:feature_importance}) and
Wilcoxon effect sizes:
co-author citation rate ($w=4.0$),
clique strength ($3.5$),
reciprocity rate ($3.5$),
outgoing HHI ($3.0$),
self-citation rate ($2.0$),
and journal endogamy rate ($2.0$).

Authors are ranked by $S_i$ in descending order.
An individual feature is flagged when its $z$-score exceeds $3\sigma$
or $5\sigma$.

\subsubsection*{Key findings}

Of the 189 outliers, 162 (85.7\%) belong to the Case (bottom-tier)
group, confirming the precision of the Isolation Forest detector.
Grouped by ASJC general field, outlier observations concentrate in
Medicine ($n=77$) and Engineering ($n=71$), followed by Social
Sciences ($n=46$) and Natural Sciences ($n=3$); note that a small
number of authors are matched in more than one field.
A total of 131 outliers (69.3\%) belong to at least one syndicate
(connected component among mutual outlier citation links), with the
largest syndicate comprising 23~members (see Figure~\ref{fig:largest-syndicate}).

Table~\ref{tab:suspicious-top10} presents the ten highest-scoring
authors.
All ten exhibit co-author citation rates exceeding $5\sigma$;
seven also show reciprocity rates beyond $5\sigma$.
The two top-ranked authors, both in agricultural engineering,
share an identical suspiciousness score ($S = 658.6$) and belong to
the same 23-member syndicate, consistent with coordinated behaviour.
The presence of authors writing in Korean, Russian, Ukrainian, and
Portuguese underscores the international scope of the phenomenon.

Table~\ref{tab:author-audit} provides a publication-level audit of
the same ten authors.
Typical portfolios are small (4--14 works, median~5), concentrated
in a single regional journal, and characterised by near-zero
outgoing citations in the global network yet elevated reciprocal
citation links, a pattern consistent with closed-loop mutual
citation within a small community of co-authors.

%% ── Table 1: Top-10 suspicious authors ────────────────────────────
\begin{table}[t]
\centering
\caption{Top-10 most suspicious authors ranked by composite
  suspiciousness score $S$ (Eq.~\ref{eq:susp}).}
\label{tab:suspicious-top10}
\scriptsize
\begin{tabularx}{\linewidth}{rl l l r r X}
\toprule
Rank & ORCID & Name & Subject & $S$ & Works & Primary anomalies \\
\midrule
  1 & \texttt{0000-0001-5794-5580} & O.\ Kochuk-Yashchenko    & Eng.  & 658.6 & 12 & co-auth.\ cit., clique str., recip.\ ($>5\sigma$) \\
  2 & \texttt{0000-0002-1998-6290} & D.\ Kucher               & Eng.  & 658.6 & 13 & co-auth.\ cit., clique str., recip.\ ($>5\sigma$) \\
  3 & \texttt{0000-0001-9044-3400} & E.\ Ivanova              & Med.  & 361.7 &  4 & co-auth.\ cit., recip., out.\ HHI ($>5\sigma$) \\
  4 & \texttt{0000-0001-5541-5180} & L.\ A.\ dos Santos Franz & Soc.  & 207.6 & 14 & co-auth.\ cit.\ ($>5\sigma$); self-cit.\ ($>3\sigma$) \\
  5 & \texttt{0000-0002-0836-9510} & J.\ Ryu                  & Eng.  & 204.9 &  4 & co-auth.\ cit., clique str., recip.\ ($>5\sigma$) \\
  6 & \texttt{0000-0002-5354-3020} & S.\ H.\ Lee              & Eng.  & 203.4 &  5 & co-auth.\ cit., clique str., recip.\ ($>5\sigma$) \\
  7 & \texttt{0000-0002-5962-4470} & M.\ L.\ Grigio           & Eng.  & 203.0 &  5 & co-auth.\ cit., recip., out.\ HHI ($>5\sigma$) \\
  8 & \texttt{0000-0002-3497-8870} & K.\ Nam                  & Eng.  & 184.7 &  4 & co-auth.\ cit., recip.\ ($>5\sigma$); endog.\ ($>5\sigma$) \\
  9 & \texttt{0000-0001-6877-9530} & D.\ R.\ Amirov           & Eng.  & 180.4 &  5 & co-auth.\ cit., recip.\ ($>5\sigma$); self-cit.\ ($>3\sigma$) \\
 10 & \texttt{0000-0003-1664-6810} & D.\ Morozova             & Eng.  & 180.4 &  4 & co-auth.\ cit., recip.\ ($>5\sigma$); self-cit.\ ($>3\sigma$) \\
\bottomrule
\end{tabularx}

\medskip
{\footnotesize\textit{Note.}
  Suspiciousness score is a weighted $z$-score composite of six
  behavioural features.
  Anomalies indicate features exceeding $3\sigma$ or $5\sigma$
  of the population mean.
  These patterns are \emph{statistically anomalous}; they do not
  constitute proof of misconduct.
  Names are initialised for brevity; full ORCIDs are verifiable via
  \texttt{https://orcid.org/[ORCID]}.}
\end{table}

%% ── Table 2: Publication audit ────────────────────────────────────
\begin{table}[t]
\centering
\caption{Publication audit profiles for top-10 suspicious authors.}
\label{tab:author-audit}
\scriptsize
\begin{tabularx}{\linewidth}{rl rrr rrX}
\toprule
Rank & ORCID & Works & Years & Out & In & Recip. & Primary journal \\
\midrule
  1 & \texttt{0000-0001-5794-5580} & 12 & 2020--2024 &  0 & 1 & 1 & Bull.\ Sumy Nat.\ Agrar.\ Univ.\ \\
  2 & \texttt{0000-0002-1998-6290} & 13 & 2020--2024 &  0 & 1 & 1 & Anim.\ Breed.\ Genet.\ \\
  3 & \texttt{0000-0001-9044-3400} &  4 & 2020--2023 &  1 & 1 & 1 & Laser Medicine \\
  4 & \texttt{0000-0001-5541-5180} & 14 & 2020--2024 &  0 & 0 & 0 & Exacta \\
  5 & \texttt{0000-0002-0836-9510} &  4 & 2020--2022 &  0 & 0 & 1 & Korean J.\ Soil Sci.\ Fert.\ \\
  6 & \texttt{0000-0002-5354-3020} &  5 & 2020--2022 &  0 & 0 & 1 & Korean J.\ Soil Sci.\ Fert.\ \\
  7 & \texttt{0000-0002-5962-4470} &  5 & 2020--2022 &  0 & 0 & 1 & Acta Sci.\ Agronomy \\
  8 & \texttt{0000-0002-3497-8870} &  4 & 2021--2023 &  0 & 0 & 1 & Korean J.\ Plant Taxon.\ \\
  9 & \texttt{0000-0001-6877-9530} &  5 & 2021--2023 &  0 & 0 & 1 & Sci.\ Notes Kazan Vet.\ Acad.\ \\
 10 & \texttt{0000-0003-1664-6810} &  4 & 2022--2023 &  0 & 0 & 1 & Int.\ J.\ Vet.\ Med.\ \\
\bottomrule
\end{tabularx}

\medskip
{\footnotesize\textit{Note.}
  Out = total outgoing citations in the directed citation network;
  In = total incoming; Recip.\ = reciprocal citation pairs.
  ``Primary journal'' is the journal in which the author published
  the most works within the 2020--2024 study window.}
\end{table}

\subsubsection*{Implications for editorial practice}

The concentration of suspicious behaviour in small, regional journals
, several of which do not appear in major indexing databases,
suggests that coordinated citation inflation is disproportionately
prevalent in venues with limited editorial oversight.
The publication-audit data show that the highest-scoring authors
maintain very small portfolios (median~5 works), yet achieve
anomalous citation metrics through closed-loop mutual referencing.
This pattern is consistent with prior work on ``citation stacking''
\citep{heneberg2016} and suggests that journal-level surveillance
of incoming self-citation rates from small author cliques could serve
as an effective early-warning signal.

%% (place after the Mixing Matrix subsection, before \section{Discussion})
%% ──────────────────────────────────────────────────────────────────

\subsection{Cross-referencing Suspicious Authors}\label{sec:suspicious-authors}

To move from aggregate outlier statistics to individual-level findings,
we cross-reference the 189 flagged outliers against real-world author
identities and publication records.
Each ORCID identifier is resolved to a name via the ORCID public
API,\footnote{\url{https://pub.orcid.org}} and publication portfolios
are reconstructed by joining author\,--\,work associations to the
Crossref works table and resolving journal titles through the ISSN
catalogue.

\subsubsection*{Suspiciousness scoring}

We define a composite \emph{suspiciousness score} as the weighted sum
of population-normalised $z$-scores across six behavioural features:
%
\begin{equation}\label{eq:susp}
  S_i = \sum_{k=1}^{6} w_k \cdot z_{ik},
  \qquad
  z_{ik} = \frac{x_{ik} - \bar{x}_k}{\sigma_k},
\end{equation}
%
where $x_{ik}$ is the raw value of feature~$k$ for author~$i$,
$\bar{x}_k$ and $\sigma_k$ are the population mean and standard
deviation, and the weights $w_k$ are chosen to reflect both
Random-Forest feature importance (Figure~\ref{fig:feature_importance}) and
Wilcoxon effect sizes:
co-author citation rate ($w=4.0$),
clique strength ($3.5$),
reciprocity rate ($3.5$),
outgoing HHI ($3.0$),
self-citation rate ($2.0$),
and journal endogamy rate ($2.0$).

Authors are ranked by $S_i$ in descending order.
An individual feature is flagged when its $z$-score exceeds $3\sigma$
or $5\sigma$.

\subsubsection*{Key findings}

Of the 189 outliers, 162 (85.7\%) belong to the Case (bottom-tier)
group, confirming the precision of the Isolation Forest detector.
Grouped by ASJC general field, outlier observations concentrate in
Medicine ($n=77$) and Engineering ($n=71$), followed by Social
Sciences ($n=46$) and Natural Sciences ($n=3$); note that a small
number of authors are matched in more than one field.
A total of 131 outliers (69.3\%) belong to at least one syndicate
(connected component among mutual outlier citation links), with the
largest syndicate comprising 23~members (see Figure~\ref{fig:largest-syndicate}).

Table~\ref{tab:suspicious-top10} presents the ten highest-scoring
authors.
All ten exhibit co-author citation rates exceeding $5\sigma$;
seven also show reciprocity rates beyond $5\sigma$.
The two top-ranked authors, both in agricultural engineering,
share an identical suspiciousness score ($S = 658.6$) and belong to
the same 23-member syndicate, consistent with coordinated behaviour.
The presence of authors writing in Korean, Russian, Ukrainian, and
Portuguese underscores the international scope of the phenomenon.

Table~\ref{tab:author-audit} provides a publication-level audit of
the same ten authors.
Typical portfolios are small (4--14 works, median~5), concentrated
in a single regional journal, and characterised by near-zero
outgoing citations in the global network yet elevated reciprocal
citation links, a pattern consistent with closed-loop mutual
citation within a small community of co-authors.

%% ── Table 1: Top-10 suspicious authors ────────────────────────────
\begin{table}[t]
\centering
\caption{Top-10 most suspicious authors ranked by composite
  suspiciousness score $S$ (Eq.~\ref{eq:susp}).}
\label{tab:suspicious-top10}
\scriptsize
\begin{tabularx}{\linewidth}{rl l l r r X}
\toprule
Rank & ORCID & Name & Subject & $S$ & Works & Primary anomalies \\
\midrule
  1 & \texttt{0000-0001-5794-5580} & O.\ Kochuk-Yashchenko    & Eng.  & 658.6 & 12 & co-auth.\ cit., clique str., recip.\ ($>5\sigma$) \\
  2 & \texttt{0000-0002-1998-6290} & D.\ Kucher               & Eng.  & 658.6 & 13 & co-auth.\ cit., clique str., recip.\ ($>5\sigma$) \\
  3 & \texttt{0000-0001-9044-3400} & E.\ Ivanova              & Med.  & 361.7 &  4 & co-auth.\ cit., recip., out.\ HHI ($>5\sigma$) \\
  4 & \texttt{0000-0001-5541-5180} & L.\ A.\ dos Santos Franz & Soc.  & 207.6 & 14 & co-auth.\ cit.\ ($>5\sigma$); self-cit.\ ($>3\sigma$) \\
  5 & \texttt{0000-0002-0836-9510} & J.\ Ryu                  & Eng.  & 204.9 &  4 & co-auth.\ cit., clique str., recip.\ ($>5\sigma$) \\
  6 & \texttt{0000-0002-5354-3020} & S.\ H.\ Lee              & Eng.  & 203.4 &  5 & co-auth.\ cit., clique str., recip.\ ($>5\sigma$) \\
  7 & \texttt{0000-0002-5962-4470} & M.\ L.\ Grigio           & Eng.  & 203.0 &  5 & co-auth.\ cit., recip., out.\ HHI ($>5\sigma$) \\
  8 & \texttt{0000-0002-3497-8870} & K.\ Nam                  & Eng.  & 184.7 &  4 & co-auth.\ cit., recip.\ ($>5\sigma$); endog.\ ($>5\sigma$) \\
  9 & \texttt{0000-0001-6877-9530} & D.\ R.\ Amirov           & Eng.  & 180.4 &  5 & co-auth.\ cit., recip.\ ($>5\sigma$); self-cit.\ ($>3\sigma$) \\
 10 & \texttt{0000-0003-1664-6810} & D.\ Morozova             & Eng.  & 180.4 &  4 & co-auth.\ cit., recip.\ ($>5\sigma$); self-cit.\ ($>3\sigma$) \\
\bottomrule
\end{tabularx}

\medskip
{\footnotesize\textit{Note.}
  Suspiciousness score is a weighted $z$-score composite of six
  behavioural features.
  Anomalies indicate features exceeding $3\sigma$ or $5\sigma$
  of the population mean.
  These patterns are \emph{statistically anomalous}; they do not
  constitute proof of misconduct.
  Names are initialised for brevity; full ORCIDs are verifiable via
  \texttt{https://orcid.org/[ORCID]}.}
\end{table}

%% ── Table 2: Publication audit ────────────────────────────────────
\begin{table}[t]
\centering
\caption{Publication audit profiles for top-10 suspicious authors.}
\label{tab:author-audit}
\scriptsize
\begin{tabularx}{\linewidth}{rl rrr rrX}
\toprule
Rank & ORCID & Works & Years & Out & In & Recip. & Primary journal \\
\midrule
  1 & \texttt{0000-0001-5794-5580} & 12 & 2020--2024 &  0 & 1 & 1 & Bull.\ Sumy Nat.\ Agrar.\ Univ.\ \\
  2 & \texttt{0000-0002-1998-6290} & 13 & 2020--2024 &  0 & 1 & 1 & Anim.\ Breed.\ Genet.\ \\
  3 & \texttt{0000-0001-9044-3400} &  4 & 2020--2023 &  1 & 1 & 1 & Laser Medicine \\
  4 & \texttt{0000-0001-5541-5180} & 14 & 2020--2024 &  0 & 0 & 0 & Exacta \\
  5 & \texttt{0000-0002-0836-9510} &  4 & 2020--2022 &  0 & 0 & 1 & Korean J.\ Soil Sci.\ Fert.\ \\
  6 & \texttt{0000-0002-5354-3020} &  5 & 2020--2022 &  0 & 0 & 1 & Korean J.\ Soil Sci.\ Fert.\ \\
  7 & \texttt{0000-0002-5962-4470} &  5 & 2020--2022 &  0 & 0 & 1 & Acta Sci.\ Agronomy \\
  8 & \texttt{0000-0002-3497-8870} &  4 & 2021--2023 &  0 & 0 & 1 & Korean J.\ Plant Taxon.\ \\
  9 & \texttt{0000-0001-6877-9530} &  5 & 2021--2023 &  0 & 0 & 1 & Sci.\ Notes Kazan Vet.\ Acad.\ \\
 10 & \texttt{0000-0003-1664-6810} &  4 & 2022--2023 &  0 & 0 & 1 & Int.\ J.\ Vet.\ Med.\ \\
\bottomrule
\end{tabularx}

\medskip
{\footnotesize\textit{Note.}
  Out = total outgoing citations in the directed citation network;
  In = total incoming; Recip.\ = reciprocal citation pairs.
  ``Primary journal'' is the journal in which the author published
  the most works within the 2020--2024 study window.}
\end{table}

\subsubsection*{Implications for editorial practice}

The concentration of suspicious behaviour in small, regional journals
, several of which do not appear in major indexing databases,
suggests that coordinated citation inflation is disproportionately
prevalent in venues with limited editorial oversight.
The publication-audit data show that the highest-scoring authors
maintain very small portfolios (median~5 works), yet achieve
anomalous citation metrics through closed-loop mutual referencing.
This pattern is consistent with prior work on ``citation stacking''
\citep{heneberg2016} and suggests that journal-level surveillance
of incoming self-citation rates from small author cliques could serve
as an effective early-warning signal.

%% (place after the Mixing Matrix subsection, before \section{Discussion})
%% ──────────────────────────────────────────────────────────────────

\subsection{Cross-referencing Suspicious Authors}\label{sec:suspicious-authors}

To move from aggregate outlier statistics to individual-level findings,
we cross-reference the 189 flagged outliers against real-world author
identities and publication records.
Each ORCID identifier is resolved to a name via the ORCID public
API,\footnote{\url{https://pub.orcid.org}} and publication portfolios
are reconstructed by joining author\,--\,work associations to the
Crossref works table and resolving journal titles through the ISSN
catalogue.

\subsubsection*{Suspiciousness scoring}

We define a composite \emph{suspiciousness score} as the weighted sum
of population-normalised $z$-scores across six behavioural features:
%
\begin{equation}\label{eq:susp}
  S_i = \sum_{k=1}^{6} w_k \cdot z_{ik},
  \qquad
  z_{ik} = \frac{x_{ik} - \bar{x}_k}{\sigma_k},
\end{equation}
%
where $x_{ik}$ is the raw value of feature~$k$ for author~$i$,
$\bar{x}_k$ and $\sigma_k$ are the population mean and standard
deviation, and the weights $w_k$ are chosen to reflect both
Random-Forest feature importance (Figure~\ref{fig:feature_importance}) and
Wilcoxon effect sizes:
co-author citation rate ($w=4.0$),
clique strength ($3.5$),
reciprocity rate ($3.5$),
outgoing HHI ($3.0$),
self-citation rate ($2.0$),
and journal endogamy rate ($2.0$).

Authors are ranked by $S_i$ in descending order.
An individual feature is flagged when its $z$-score exceeds $3\sigma$
or $5\sigma$.

\subsubsection*{Key findings}

Of the 189 outliers, 162 (85.7\%) belong to the Case (bottom-tier)
group, confirming the precision of the Isolation Forest detector.
Grouped by ASJC general field, outlier observations concentrate in
Medicine ($n=77$) and Engineering ($n=71$), followed by Social
Sciences ($n=46$) and Natural Sciences ($n=3$); note that a small
number of authors are matched in more than one field.
A total of 131 outliers (69.3\%) belong to at least one syndicate
(connected component among mutual outlier citation links), with the
largest syndicate comprising 23~members (see Figure~\ref{fig:largest-syndicate}).

Table~\ref{tab:suspicious-top10} presents the ten highest-scoring
authors.
All ten exhibit co-author citation rates exceeding $5\sigma$;
seven also show reciprocity rates beyond $5\sigma$.
The two top-ranked authors, both in agricultural engineering,
share an identical suspiciousness score ($S = 658.6$) and belong to
the same 23-member syndicate, consistent with coordinated behaviour.
The presence of authors writing in Korean, Russian, Ukrainian, and
Portuguese underscores the international scope of the phenomenon.

Table~\ref{tab:author-audit} provides a publication-level audit of
the same ten authors.
Typical portfolios are small (4--14 works, median~5), concentrated
in a single regional journal, and characterised by near-zero
outgoing citations in the global network yet elevated reciprocal
citation links, a pattern consistent with closed-loop mutual
citation within a small community of co-authors.

%% ── Table 1: Top-10 suspicious authors ────────────────────────────
\begin{table}[t]
\centering
\caption{Top-10 most suspicious authors ranked by composite
  suspiciousness score $S$ (Eq.~\ref{eq:susp}).}
\label{tab:suspicious-top10}
\scriptsize
\begin{tabularx}{\linewidth}{rl l l r r X}
\toprule
Rank & ORCID & Name & Subject & $S$ & Works & Primary anomalies \\
\midrule
  1 & \texttt{0000-0001-5794-5580} & O.\ Kochuk-Yashchenko    & Eng.  & 658.6 & 12 & co-auth.\ cit., clique str., recip.\ ($>5\sigma$) \\
  2 & \texttt{0000-0002-1998-6290} & D.\ Kucher               & Eng.  & 658.6 & 13 & co-auth.\ cit., clique str., recip.\ ($>5\sigma$) \\
  3 & \texttt{0000-0001-9044-3400} & E.\ Ivanova              & Med.  & 361.7 &  4 & co-auth.\ cit., recip., out.\ HHI ($>5\sigma$) \\
  4 & \texttt{0000-0001-5541-5180} & L.\ A.\ dos Santos Franz & Soc.  & 207.6 & 14 & co-auth.\ cit.\ ($>5\sigma$); self-cit.\ ($>3\sigma$) \\
  5 & \texttt{0000-0002-0836-9510} & J.\ Ryu                  & Eng.  & 204.9 &  4 & co-auth.\ cit., clique str., recip.\ ($>5\sigma$) \\
  6 & \texttt{0000-0002-5354-3020} & S.\ H.\ Lee              & Eng.  & 203.4 &  5 & co-auth.\ cit., clique str., recip.\ ($>5\sigma$) \\
  7 & \texttt{0000-0002-5962-4470} & M.\ L.\ Grigio           & Eng.  & 203.0 &  5 & co-auth.\ cit., recip., out.\ HHI ($>5\sigma$) \\
  8 & \texttt{0000-0002-3497-8870} & K.\ Nam                  & Eng.  & 184.7 &  4 & co-auth.\ cit., recip.\ ($>5\sigma$); endog.\ ($>5\sigma$) \\
  9 & \texttt{0000-0001-6877-9530} & D.\ R.\ Amirov           & Eng.  & 180.4 &  5 & co-auth.\ cit., recip.\ ($>5\sigma$); self-cit.\ ($>3\sigma$) \\
 10 & \texttt{0000-0003-1664-6810} & D.\ Morozova             & Eng.  & 180.4 &  4 & co-auth.\ cit., recip.\ ($>5\sigma$); self-cit.\ ($>3\sigma$) \\
\bottomrule
\end{tabularx}

\medskip
{\footnotesize\textit{Note.}
  Suspiciousness score is a weighted $z$-score composite of six
  behavioural features.
  Anomalies indicate features exceeding $3\sigma$ or $5\sigma$
  of the population mean.
  These patterns are \emph{statistically anomalous}; they do not
  constitute proof of misconduct.
  Names are initialised for brevity; full ORCIDs are verifiable via
  \texttt{https://orcid.org/[ORCID]}.}
\end{table}

%% ── Table 2: Publication audit ────────────────────────────────────
\begin{table}[t]
\centering
\caption{Publication audit profiles for top-10 suspicious authors.}
\label{tab:author-audit}
\scriptsize
\begin{tabularx}{\linewidth}{rl rrr rrX}
\toprule
Rank & ORCID & Works & Years & Out & In & Recip. & Primary journal \\
\midrule
  1 & \texttt{0000-0001-5794-5580} & 12 & 2020--2024 &  0 & 1 & 1 & Bull.\ Sumy Nat.\ Agrar.\ Univ.\ \\
  2 & \texttt{0000-0002-1998-6290} & 13 & 2020--2024 &  0 & 1 & 1 & Anim.\ Breed.\ Genet.\ \\
  3 & \texttt{0000-0001-9044-3400} &  4 & 2020--2023 &  1 & 1 & 1 & Laser Medicine \\
  4 & \texttt{0000-0001-5541-5180} & 14 & 2020--2024 &  0 & 0 & 0 & Exacta \\
  5 & \texttt{0000-0002-0836-9510} &  4 & 2020--2022 &  0 & 0 & 1 & Korean J.\ Soil Sci.\ Fert.\ \\
  6 & \texttt{0000-0002-5354-3020} &  5 & 2020--2022 &  0 & 0 & 1 & Korean J.\ Soil Sci.\ Fert.\ \\
  7 & \texttt{0000-0002-5962-4470} &  5 & 2020--2022 &  0 & 0 & 1 & Acta Sci.\ Agronomy \\
  8 & \texttt{0000-0002-3497-8870} &  4 & 2021--2023 &  0 & 0 & 1 & Korean J.\ Plant Taxon.\ \\
  9 & \texttt{0000-0001-6877-9530} &  5 & 2021--2023 &  0 & 0 & 1 & Sci.\ Notes Kazan Vet.\ Acad.\ \\
 10 & \texttt{0000-0003-1664-6810} &  4 & 2022--2023 &  0 & 0 & 1 & Int.\ J.\ Vet.\ Med.\ \\
\bottomrule
\end{tabularx}

\medskip
{\footnotesize\textit{Note.}
  Out = total outgoing citations in the directed citation network;
  In = total incoming; Recip.\ = reciprocal citation pairs.
  ``Primary journal'' is the journal in which the author published
  the most works within the 2020--2024 study window.}
\end{table}

\subsubsection*{Implications for editorial practice}

The concentration of suspicious behaviour in small, regional journals
, several of which do not appear in major indexing databases,
suggests that coordinated citation inflation is disproportionately
prevalent in venues with limited editorial oversight.
The publication-audit data show that the highest-scoring authors
maintain very small portfolios (median~5 works), yet achieve
anomalous citation metrics through closed-loop mutual referencing.
This pattern is consistent with prior work on ``citation stacking''
\citep{heneberg2016} and suggests that journal-level surveillance
of incoming self-citation rates from small author cliques could serve
as an effective early-warning signal.

%% (place after the Mixing Matrix subsection, before \section{Discussion})
%% ──────────────────────────────────────────────────────────────────

\subsection{Cross-referencing Suspicious Authors}\label{sec:suspicious-authors}

To move from aggregate outlier statistics to individual-level findings,
we cross-reference the 189 flagged outliers against real-world author
identities and publication records.
Each ORCID identifier is resolved to a name via the ORCID public
API,\footnote{\url{https://pub.orcid.org}} and publication portfolios
are reconstructed by joining author\,--\,work associations to the
Crossref works table and resolving journal titles through the ISSN
catalogue.

\subsubsection*{Suspiciousness scoring}

We define a composite \emph{suspiciousness score} as the weighted sum
of population-normalised $z$-scores across six behavioural features:
%
\begin{equation}\label{eq:susp}
  S_i = \sum_{k=1}^{6} w_k \cdot z_{ik},
  \qquad
  z_{ik} = \frac{x_{ik} - \bar{x}_k}{\sigma_k},
\end{equation}
%
where $x_{ik}$ is the raw value of feature~$k$ for author~$i$,
$\bar{x}_k$ and $\sigma_k$ are the population mean and standard
deviation, and the weights $w_k$ are chosen to reflect both
Random-Forest feature importance (Figure~\ref{fig:feature_importance}) and
Wilcoxon effect sizes:
co-author citation rate ($w=4.0$),
clique strength ($3.5$),
reciprocity rate ($3.5$),
outgoing HHI ($3.0$),
self-citation rate ($2.0$),
and journal endogamy rate ($2.0$).

Authors are ranked by $S_i$ in descending order.
An individual feature is flagged when its $z$-score exceeds $3\sigma$
or $5\sigma$.

\subsubsection*{Key findings}

Of the 189 outliers, 162 (85.7\%) belong to the Case (bottom-tier)
group, confirming the precision of the Isolation Forest detector.
Grouped by ASJC general field, outlier observations concentrate in
Medicine ($n=77$) and Engineering ($n=71$), followed by Social
Sciences ($n=46$) and Natural Sciences ($n=3$); note that a small
number of authors are matched in more than one field.
A total of 131 outliers (69.3\%) belong to at least one syndicate
(connected component among mutual outlier citation links), with the
largest syndicate comprising 23~members (see Figure~\ref{fig:largest-syndicate}).

Table~\ref{tab:suspicious-top10} presents the ten highest-scoring
authors.
All ten exhibit co-author citation rates exceeding $5\sigma$;
seven also show reciprocity rates beyond $5\sigma$.
The two top-ranked authors, both in agricultural engineering,
share an identical suspiciousness score ($S = 658.6$) and belong to
the same 23-member syndicate, consistent with coordinated behaviour.
The presence of authors writing in Korean, Russian, Ukrainian, and
Portuguese underscores the international scope of the phenomenon.

Table~\ref{tab:author-audit} provides a publication-level audit of
the same ten authors.
Typical portfolios are small (4--14 works, median~5), concentrated
in a single regional journal, and characterised by near-zero
outgoing citations in the global network yet elevated reciprocal
citation links, a pattern consistent with closed-loop mutual
citation within a small community of co-authors.

%% ── Table 1: Top-10 suspicious authors ────────────────────────────
\begin{table}[t]
\centering
\caption{Top-10 most suspicious authors ranked by composite
  suspiciousness score $S$ (Eq.~\ref{eq:susp}).}
\label{tab:suspicious-top10}
\scriptsize
\begin{tabularx}{\linewidth}{rl l l r r X}
\toprule
Rank & ORCID & Name & Subject & $S$ & Works & Primary anomalies \\
\midrule
  1 & \texttt{0000-0001-5794-5580} & O.\ Kochuk-Yashchenko    & Eng.  & 658.6 & 12 & co-auth.\ cit., clique str., recip.\ ($>5\sigma$) \\
  2 & \texttt{0000-0002-1998-6290} & D.\ Kucher               & Eng.  & 658.6 & 13 & co-auth.\ cit., clique str., recip.\ ($>5\sigma$) \\
  3 & \texttt{0000-0001-9044-3400} & E.\ Ivanova              & Med.  & 361.7 &  4 & co-auth.\ cit., recip., out.\ HHI ($>5\sigma$) \\
  4 & \texttt{0000-0001-5541-5180} & L.\ A.\ dos Santos Franz & Soc.  & 207.6 & 14 & co-auth.\ cit.\ ($>5\sigma$); self-cit.\ ($>3\sigma$) \\
  5 & \texttt{0000-0002-0836-9510} & J.\ Ryu                  & Eng.  & 204.9 &  4 & co-auth.\ cit., clique str., recip.\ ($>5\sigma$) \\
  6 & \texttt{0000-0002-5354-3020} & S.\ H.\ Lee              & Eng.  & 203.4 &  5 & co-auth.\ cit., clique str., recip.\ ($>5\sigma$) \\
  7 & \texttt{0000-0002-5962-4470} & M.\ L.\ Grigio           & Eng.  & 203.0 &  5 & co-auth.\ cit., recip., out.\ HHI ($>5\sigma$) \\
  8 & \texttt{0000-0002-3497-8870} & K.\ Nam                  & Eng.  & 184.7 &  4 & co-auth.\ cit., recip.\ ($>5\sigma$); endog.\ ($>5\sigma$) \\
  9 & \texttt{0000-0001-6877-9530} & D.\ R.\ Amirov           & Eng.  & 180.4 &  5 & co-auth.\ cit., recip.\ ($>5\sigma$); self-cit.\ ($>3\sigma$) \\
 10 & \texttt{0000-0003-1664-6810} & D.\ Morozova             & Eng.  & 180.4 &  4 & co-auth.\ cit., recip.\ ($>5\sigma$); self-cit.\ ($>3\sigma$) \\
\bottomrule
\end{tabularx}

\medskip
{\footnotesize\textit{Note.}
  Suspiciousness score is a weighted $z$-score composite of six
  behavioural features.
  Anomalies indicate features exceeding $3\sigma$ or $5\sigma$
  of the population mean.
  These patterns are \emph{statistically anomalous}; they do not
  constitute proof of misconduct.
  Names are initialised for brevity; full ORCIDs are verifiable via
  \texttt{https://orcid.org/[ORCID]}.}
\end{table}

%% ── Table 2: Publication audit ────────────────────────────────────
\begin{table}[t]
\centering
\caption{Publication audit profiles for top-10 suspicious authors.}
\label{tab:author-audit}
\scriptsize
\begin{tabularx}{\linewidth}{rl rrr rrX}
\toprule
Rank & ORCID & Works & Years & Out & In & Recip. & Primary journal \\
\midrule
  1 & \texttt{0000-0001-5794-5580} & 12 & 2020--2024 &  0 & 1 & 1 & Bull.\ Sumy Nat.\ Agrar.\ Univ.\ \\
  2 & \texttt{0000-0002-1998-6290} & 13 & 2020--2024 &  0 & 1 & 1 & Anim.\ Breed.\ Genet.\ \\
  3 & \texttt{0000-0001-9044-3400} &  4 & 2020--2023 &  1 & 1 & 1 & Laser Medicine \\
  4 & \texttt{0000-0001-5541-5180} & 14 & 2020--2024 &  0 & 0 & 0 & Exacta \\
  5 & \texttt{0000-0002-0836-9510} &  4 & 2020--2022 &  0 & 0 & 1 & Korean J.\ Soil Sci.\ Fert.\ \\
  6 & \texttt{0000-0002-5354-3020} &  5 & 2020--2022 &  0 & 0 & 1 & Korean J.\ Soil Sci.\ Fert.\ \\
  7 & \texttt{0000-0002-5962-4470} &  5 & 2020--2022 &  0 & 0 & 1 & Acta Sci.\ Agronomy \\
  8 & \texttt{0000-0002-3497-8870} &  4 & 2021--2023 &  0 & 0 & 1 & Korean J.\ Plant Taxon.\ \\
  9 & \texttt{0000-0001-6877-9530} &  5 & 2021--2023 &  0 & 0 & 1 & Sci.\ Notes Kazan Vet.\ Acad.\ \\
 10 & \texttt{0000-0003-1664-6810} &  4 & 2022--2023 &  0 & 0 & 1 & Int.\ J.\ Vet.\ Med.\ \\
\bottomrule
\end{tabularx}

\medskip
{\footnotesize\textit{Note.}
  Out = total outgoing citations in the directed citation network;
  In = total incoming; Recip.\ = reciprocal citation pairs.
  ``Primary journal'' is the journal in which the author published
  the most works within the 2020--2024 study window.}
\end{table}

\subsubsection*{Implications for editorial practice}

The concentration of suspicious behaviour in small, regional journals
, several of which do not appear in major indexing databases,
suggests that coordinated citation inflation is disproportionately
prevalent in venues with limited editorial oversight.
The publication-audit data show that the highest-scoring authors
maintain very small portfolios (median~5 works), yet achieve
anomalous citation metrics through closed-loop mutual referencing.
This pattern is consistent with prior work on ``citation stacking''
\citep{heneberg2016} and suggests that journal-level surveillance
of incoming self-citation rates from small author cliques could serve
as an effective early-warning signal.
